\subsection{Оценка цикла по формуле Уэбстера}

Для варианта 2 входной спрос на каждой полосе $q$ равен $0{,}95$,
делённому на средний интервал прибытия на въезд и умноженному
на долю соответствующего манёвра. В исходном графе каждый манёвр
закреплён за своей полосой.

Поток насыщения измеряли на базовом плане (42 с, по 18 с зелёного)
при шестикратном базовом входном спросе. Для каждого SEED 201–203
отбросили 8 циклов и использовали следующие 20. В начале зелёного
отбирали две остановившиеся машины в голове очереди на каждой полосе.
Проезд стоп-линии фиксировали при переходе на траекторию перекрёстка.
Для каждой полосы число интервалов между проездами таких пар делили
на их суммарную длительность. Так получили поток насыщения $s$ каждой
полосы. Способ исключает интервалы зелёного после исчерпания очереди.

Все полосы фазы получают одинаковое время зелёного. Его требуемую долю
определяет полоса с максимальным отношением $q/s$. На каждой фазе
критическими оказались полосы движения прямо, на которые приходится 60\% спроса.

Потерянное время $L=2\cdot3=6$ с приходится на два интервала общего красного.
Модель сразу разрешает движение с заданной скоростью и не задаёт отдельной
задержки реакции или разгона.

\begin{center}
\begin{tabular}{llrrr}
Фаза & Критические полосы & $q_i$, маш/с & $s_i$, маш/с & $y_i=q_i/s_i$ \\
NS & N1, S1 & 0,2375 & 1,4286 & 0,1663 \\
EW & E1, W1 & 0,1781 & 1,4286 & 0,1247 \\
\end{tabular}
\end{center}

При $Y=\sum_i y_i=0{,}2909$ получаем
$C_0=(1{,}5L+5)/(1-Y)=19{,}74$ с. Это на 22,26 с короче базового
цикла и на 10,26 с ниже нижней границы сетки поиска (30 с).

Дискретная сетка, ограничения фаз и распределение машин по полосам
влияют на план, который выбирается в симуляции. Критерий поиска по T-41
фиксируется отдельно до перебора сетки.
Оценка предварительная. После T-11 повторить замер на независимых потоках.
После T-45 подставить выбранный цикл и указать фактическое расхождение.

Команда из корня репозитория:
\begin{verbatim}
uv run python -m scripts.webster > results/webster.csv
\end{verbatim}
